En esta sección se describen los requisitos de gestión de información (datos) que el sistema debe gestionar. Estos requisitos especifican los tipos de información que el sistema debe almacenar, consultar y mantener a lo largo de su funcionamiento.

A diferencia de los requisitos funcionales, que describen los servicios que el sistema proporciona, los requisitos de información definen las estructuras de datos necesarias para soportar dichos servicios. La definición detallada de las relaciones entre estos datos y de sus restricciones de integridad se abordará posteriormente en el modelo conceptual de datos.

Con el objetivo de mantener la coherencia con la organización del sistema, los requisitos de información se presentan agrupados por módulos funcionales.

% ------------------------------------------------------------------------------
\subsection{M1. Gestión de acceso, alta y autorización de usuarios}
% ------------------------------------------------------------------------------

Este módulo gestiona la información necesaria para identificar a los usuarios del sistema, administrar sus permisos de acceso y registrar las actividades
relacionadas con la autenticación y la gestión de cuentas.

\subsubsection*{Usuario}

El sistema debe almacenar información básica de cada usuario registrado en la plataforma. Esta información permite identificar al usuario, gestionar su acceso al sistema y mantener información básica de contacto.

Los datos asociados a un usuario incluyen:

\begin{itemize}
    \item Identificador del usuario.
    \item Dirección de correo electrónico.
    \item Contraseña cifrada.
    \item Nombre completo.
    \item Teléfono de contacto.
    \item Empresa o entidad a la que pertenece.
    \item Rol asignado dentro del sistema.
    \item Estado de la cuenta de usuario.
    \item Imagen de perfil.
    \item Fecha del último acceso al sistema.
    \item Fecha de creación de la cuenta.
    \item Fecha de última actualización de la cuenta.
\end{itemize}

\subsubsection*{Rol}

El sistema debe gestionar los diferentes roles que determinan los permisos y
responsabilidades de los usuarios dentro de la aplicación.

Los datos asociados a un rol incluyen:

\begin{itemize}
    \item Identificador del rol.
    \item Código del rol.
    \item Nombre del rol.
    \item Descripción del rol.
\end{itemize}

\subsubsection*{Solicitud de registro}

El sistema debe almacenar las solicitudes de registro realizadas por usuarios
que desean acceder a la plataforma. Estas solicitudes podrán ser revisadas por
un administrador antes de conceder o denegar el acceso.

Los datos asociados a una solicitud de registro incluyen:

\begin{itemize}
    \item Identificador de la solicitud.
    \item Correo electrónico del solicitante.
    \item Contraseña cifrada propuesta.
    \item Nombre del solicitante.
    \item Teléfono de contacto.
    \item Empresa asociada.
    \item Rol solicitado.
    \item Estado de la solicitud.
    \item Origen de la solicitud de registro.
    \item Fecha de realización de la solicitud.
    \item Fecha de revisión de la solicitud.
    \item Usuario que realizó la revisión.
    \item Motivo de rechazo (en caso de que la solicitud sea denegada).
    \item Usuario creado a partir de la solicitud, en caso de aprobación.
\end{itemize}

\subsubsection*{Registro de administración de usuarios}

El sistema debe mantener un historial de las acciones administrativas realizadas sobre las cuentas de usuario con el objetivo de disponer de trazabilidad sobre los cambios efectuados en el sistema.

Los datos asociados a cada registro de administración incluyen:

\begin{itemize}
    \item Identificador del registro.
    \item Usuario sobre el que se realiza la acción.
    \item Usuario que realiza la acción administrativa.
    \item Tipo de acción realizada.
    \item Estado previo del usuario.
    \item Nuevo estado del usuario.
    \item Rol previo del usuario.
    \item Nuevo rol del usuario.
    \item Motivo de la acción administrativa.
    \item Notas adicionales asociadas a la acción.
    \item Fecha en la que se realizó la acción.
\end{itemize}

\subsubsection*{Registro de acceso}

El sistema debe registrar los eventos relacionados con el acceso al sistema con el objetivo de mejorar la seguridad, detectar accesos indebidos y disponer
de información para auditoría.

Los datos asociados a un registro de acceso incluyen:

\begin{itemize}
    \item Identificador del registro de acceso.
    \item Usuario asociado al evento de acceso.
    \item Dirección de correo electrónico utilizada en el intento de acceso.
    \item Tipo de evento de acceso.
    \item Resultado del intento de acceso.
    \item Motivo del fallo de acceso, en caso de error.
    \item Identificador de sesión asociado al acceso.
    \item Dirección IP desde la que se realizó el acceso.
    \item Información del agente de usuario (navegador o dispositivo).
    \item Fecha y hora del evento de acceso.
    \item Fecha de revocación de la sesión, en caso de existir.
    \item Fecha de expiración de la sesión.
\end{itemize}

% ------------------------------------------------------------------------------
\subsection{M2. Gestión comercial y clientes}
% ------------------------------------------------------------------------------

Este módulo gestiona la información relacionada con los clientes del sistema y con la actividad comercial desarrollada por los agentes comerciales. Su objetivo es centralizar la información relevante sobre las peluquerías clientes, facilitar el trabajo de los agentes comerciales durante sus visitas, permitir el seguimiento de la actividad comercial realizada y dar soporte a la planificación de rutas comerciales mediante la consulta de la información de contacto y localización de los clientes.

\subsubsection*{Cliente}

El sistema debe almacenar la información básica de cada cliente profesional atendido por el distribuidor. Esta información permite identificar cada establecimiento, registrar los datos comerciales relevantes, mantener la relación entre la peluquería como entidad de negocio y su cuenta de acceso en la plataforma, y facilitar al agente comercial la consulta de los datos de contacto y localización necesarios para la planificación de visitas y rutas comerciales.

Los datos asociados a un cliente incluyen:

\begin{itemize}
    \item Identificador del cliente.
    \item Nombre del establecimiento o peluquería.
    \item Nombre de la persona de contacto.
    \item Identificador fiscal del cliente, cuando corresponda.
    \item Dirección del establecimiento.
    \item Localidad o ciudad.
    \item Código postal.
    \item Provincia o zona.
    \item Usuario comercial responsable del cliente.
    \item Usuario asociado al cliente dentro del sistema.
    \item Observaciones o notas comerciales.
    \item Fecha de creación del registro del cliente.
    \item Fecha de última actualización de la información del cliente.
\end{itemize}

Esta información podrá ser utilizada por el sistema para mostrar al comercial los datos necesarios para localizar al cliente durante la planificación y ejecución de rutas comerciales.

\subsubsection*{Visita comercial}

El sistema debe permitir registrar las visitas comerciales realizadas por los agentes a los clientes profesionales. Estas visitas permiten documentar la actividad comercial y registrar observaciones o resultados obtenidos durante la interacción con el cliente.

Los datos asociados a una visita comercial incluyen:

\begin{itemize}
    \item Identificador de la visita.
    \item Cliente visitado.
    \item Usuario comercial que realiza la visita.
    \item Fecha y hora de la visita.
    \item Estado de la visita (planificada, realizada o cancelada).
    \item Observaciones registradas durante la visita.
    \item Resultado o conclusión de la visita.
\end{itemize}

Esta información podrá ser utilizada por el sistema para mostrar al comercial los datos necesarios para localizar al cliente durante la planificación y ejecución de rutas comerciales.

\subsubsection*{Ruta comercial}

El sistema debe permitir organizar las rutas de trabajo de los agentes comerciales. Una ruta representa la planificación de visitas que un comercial debe realizar durante una jornada o recorrido determinado. Para ello, el sistema debe permitir consultar la información de contacto y localización de los clientes incluidos en la ruta y servir de apoyo a la generación de indicaciones de navegación mediante servicios externos.

Los datos asociados a una ruta comercial incluyen:

\begin{itemize}
    \item Identificador de la ruta.
    \item Usuario comercial responsable de la ruta.
    \item Fecha de la ruta comercial.
    \item Nombre o descripción de la ruta.
    \item Estado de la ruta.
    \item Fecha de creación de la ruta.
\end{itemize}

\subsubsection*{Visita en ruta}

El sistema debe permitir registrar las visitas que forman parte de una ruta comercial concreta. Esta información permite definir las visitas comerciales incluidas en la ruta y el orden previsto en el que deben realizarse, facilitando así la organización del recorrido comercial diario.

Los datos asociados a una visita dentro de una ruta incluyen:

\begin{itemize}
    \item Identificador del registro.
    \item Ruta comercial asociada.
    \item Visita comercial incluida en la ruta.
    \item Orden previsto de la visita dentro de la ruta.
\end{itemize}

% ------------------------------------------------------------------------------
\subsection{M3. Catálogo, productos y coloración}
% ------------------------------------------------------------------------------

Este módulo gestiona la información relacionada con el catálogo de productos profesionales comercializados por el distribuidor. Su objetivo es permitir la consulta estructurada de los productos, sus características comerciales y técnicas, así como de los recursos de apoyo y de las referencias específicas asociadas a los procesos de coloración capilar.

\subsubsection*{Producto}

El sistema debe almacenar la información básica de cada producto profesional incluido en el catálogo. Esta información permite identificar el producto, clasificarlo dentro de la oferta comercial y mostrar sus características principales a los usuarios del sistema.

Los datos asociados a un producto incluyen:

\begin{itemize}
    \item Identificador del producto.
    \item Nombre del producto.
    \item Referencia comercial del producto.
    \item Descripción del producto.
    \item Categoría o familia a la que pertenece.
    \item Línea comercial asociada.
    \item Imagen principal del producto.
    \item Información técnica o modo de aplicación.
    \item Estado del producto dentro del catálogo.
    \item Precio base del producto.
    \item Proveedor asociado al producto, cuando corresponda.
    \item Fecha de creación del registro del producto.
    \item Fecha de última actualización de la información del producto.
\end{itemize}

\subsubsection*{Categoría de producto}

El sistema debe almacenar la clasificación general de los productos del catálogo con el fin de facilitar su organización y consulta estructurada.

Los datos asociados a una categoría de producto incluyen:

\begin{itemize}
    \item Identificador de la categoría.
    \item Nombre de la categoría.
    \item Descripción de la categoría.
    \item Orden o posición de visualización.
\end{itemize}

\subsubsection*{Línea comercial}

El sistema debe almacenar la información de las líneas comerciales en las que se agrupan los productos. Esta información permite organizar el catálogo según la estructura comercial utilizada por el distribuidor.

Los datos asociados a una línea comercial incluyen:

\begin{itemize}
    \item Identificador de la línea comercial.
    \item Nombre de la línea comercial.
    \item Descripción de la línea comercial.
    \item Categoría asociada.
    \item Imagen o recurso gráfico representativo.
    \item Orden o posición de visualización.
\end{itemize}

\subsubsection*{Recurso de apoyo}

El sistema debe permitir almacenar recursos de apoyo asociados a los productos o líneas comerciales del catálogo. Estos recursos facilitan el acceso a documentación técnica, catálogos comerciales o materiales formativos.

Los datos asociados a un recurso de apoyo incluyen:

\begin{itemize}
    \item Identificador del recurso.
    \item Título o nombre del recurso.
    \item Descripción del recurso.
    \item Tipo de recurso.
    \item Ubicación o enlace al recurso.
    \item Producto asociado, en caso de existir.
    \item Línea comercial asociada, en caso de existir.
    \item Fecha de creación del recurso.
\end{itemize}

\subsubsection*{Carta de color}

El sistema debe permitir almacenar cartas de color asociadas a las líneas de coloración disponibles en el catálogo. Estas cartas permiten organizar y consultar las referencias cromáticas comercializadas.

Los datos asociados a una carta de color incluyen:

\begin{itemize}
    \item Identificador de la carta de color.
    \item Nombre de la carta de color.
    \item Descripción de la carta.
    \item Línea comercial asociada.
    \item Imagen o recurso visual principal de la carta.
    \item Fecha de creación del registro de la carta.
    \item Fecha de última actualización de la carta.
\end{itemize}

\subsubsection*{Referencia de color}

El sistema debe almacenar las referencias individuales incluidas en cada carta de color. Estas referencias permiten identificar tonos concretos según la nomenclatura comercial empleada por la marca.

Los datos asociados a una referencia de color incluyen:

\begin{itemize}
    \item Identificador de la referencia de color.
    \item Carta de color asociada.
    \item Código o nomenclatura comercial del tono.
    \item Nombre o denominación del tono.
    \item Descripción del tono.
    \item Imagen o muestra visual asociada.
    \item Orden o posición dentro de la carta.
\end{itemize}

% ------------------------------------------------------------------------------
\subsection{M4. Pedidos, entregas y cobros}
% ------------------------------------------------------------------------------

Este módulo gestiona la información relacionada con el proceso comercial de venta de productos, desde la creación de pedidos hasta el seguimiento de su entrega y el registro de los cobros asociados. Su objetivo es permitir un control estructurado de las operaciones comerciales realizadas con los clientes profesionales, así como del estado de los pedidos, las incidencias de entrega y los importes pendientes o abonados.

\subsubsection*{Pedido}

El sistema debe almacenar la información general asociada a cada pedido realizado por un cliente profesional o por un usuario autorizado en su nombre. Esta información permite identificar el pedido, conocer su situación dentro del proceso comercial y relacionarlo con el cliente correspondiente.

Los datos asociados a un pedido incluyen:

\begin{itemize}
    \item Identificador del pedido.
    \item Cliente asociado al pedido.
    \item Usuario que registra el pedido.
    \item Fecha de creación del pedido.
    \item Estado del pedido.
    \item Importe total del pedido.
    \item Observaciones asociadas al pedido.
    \item Fecha de última actualización del pedido.
\end{itemize}

\subsubsection*{Línea de pedido}

El sistema debe almacenar el detalle de los productos incluidos en cada pedido. Esta información permite conocer qué productos se han solicitado y en qué cantidades.

Los datos asociados a una línea de pedido incluyen:

\begin{itemize}
    \item Identificador de la línea de pedido.
    \item Pedido asociado.
    \item Producto incluido en la línea.
    \item Cantidad solicitada.
    \item Precio unitario aplicado.
    \item Descuento aplicado, en caso de existir.
    \item Importe total de la línea.
\end{itemize}

\subsubsection*{Concepto adicional de pedido}

El sistema debe permitir registrar elementos económicos adicionales asociados a un pedido que no correspondan directamente a productos del catálogo, como gastos, ajustes o servicios complementarios.

Los datos asociados incluyen:

\begin{itemize}
    \item Identificador del concepto.
    \item Pedido asociado.
    \item Descripción del concepto.
    \item Importe asociado.
\end{itemize}

\subsubsection*{Entrega}

El sistema debe almacenar la información relativa a la entrega asociada a cada pedido. Esta información permite conocer la modalidad de entrega seleccionada, su estado, las fechas previstas o reales asociadas al proceso y las incidencias producidas durante la entrega.

Los datos asociados a una entrega incluyen:

\begin{itemize}
    \item Identificador de la entrega.
    \item Pedido asociado a la entrega.
    \item Modalidad de entrega.
    \item Usuario responsable de la entrega, en caso de entrega directa.
    \item Estado de la entrega.
    \item Fecha prevista de entrega.
    \item Fecha real de entrega.
    \item Número de bultos o cajas asociados a la entrega, cuando corresponda.
    \item Importe de contrareembolso, cuando corresponda.
    \item Observaciones de la entrega.
\end{itemize}

\subsubsection*{Incidencia de entrega}

El sistema debe permitir registrar incidencias asociadas a una entrega. Estas incidencias permiten documentar problemas detectados durante el proceso de entrega, productos pendientes de servir, entregas parciales u otras situaciones relevantes para el seguimiento logístico del pedido.

Los datos asociados a una incidencia de entrega incluyen:

\begin{itemize}
    \item Identificador de la incidencia.
    \item Entrega asociada.
    \item Tipo de incidencia.
    \item Descripción de la incidencia.
    \item Fecha de registro de la incidencia.
    \item Estado de la incidencia.
\end{itemize}

\subsubsection*{Cobro}

El sistema debe permitir registrar los cobros asociados a pedidos o entregas. Esta información permite reflejar los importes abonados por los clientes y mantener el seguimiento de pagos totales o parciales.

Los datos asociados a un cobro incluyen:

\begin{itemize}
    \item Identificador del cobro.
    \item Cliente asociado al cobro.
    \item Pedido asociado al cobro, en caso de existir.
    \item Entrega asociada al cobro, en caso de existir.
    \item Fecha del cobro.
    \item Importe cobrado.
    \item Método de cobro.
    \item Observaciones asociadas al cobro.
\end{itemize}

% ------------------------------------------------------------------------------
\subsection{M5. Gestión técnica del salón y seguimiento de servicios}
% ------------------------------------------------------------------------------

El módulo de gestión técnica del salón y seguimiento de servicios gestiona la información relacionada con los clientes finales atendidos en las peluquerías, los servicios realizados y la información técnica asociada a dichos tratamientos. Su objetivo es permitir a los profesionales del salón mantener un historial estructurado de los trabajos realizados, registrar los productos utilizados, consultar la evolución técnica de cada cliente y facilitar la generación de comunicaciones personalizadas dirigidas a dichos clientes a partir de la información técnica registrada.

La información gestionada por este módulo permite documentar los servicios realizados en el salón, registrar fórmulas y productos empleados, facilitar la consulta de tratamientos previos como apoyo a futuros trabajos técnicos y apoyar la elaboración de correos personalizados basados en el historial técnico del cliente del salón.

\subsubsection*{Cliente del salón}

El sistema debe almacenar la información básica de cada cliente final atendido por una peluquería usuaria del sistema.

Para cada cliente del salón el sistema debe registrar:

\begin{itemize}
    \item Identificador del cliente del salón.
    \item Cliente profesional o peluquería a la que pertenece.
    \item Nombre del cliente del salón.
    \item Teléfono de contacto, en caso de registrarse.
    \item Correo electrónico de contacto, en caso de registrarse.
    \item Observaciones generales asociadas al cliente.
    \item Fecha de creación del registro.
    \item Fecha de última actualización de la información.
\end{itemize}

\subsubsection*{Correo generado a partir de la ficha técnica}

El sistema debe permitir generar correos electrónicos dirigidos a los clientes del salón a partir de la información registrada en su ficha técnica. Estos correos podrán incluir información relevante sobre tratamientos realizados, productos utilizados, sugerencias basadas en el historial técnico del cliente y promociones que el profesional del salón desee comunicar.

Para cada correo generado el sistema debe poder gestionar:

\begin{itemize}
    \item Cliente del salón destinatario del correo.
    \item Dirección de correo electrónico del destinatario.
    \item Asunto del correo.
    \item Contenido generado a partir de la información técnica registrada.
    \item Contenido personalizado por el profesional antes del envío.
    \item Fecha y hora de generación del correo, en caso de registrarse.
\end{itemize}

\subsubsection*{Servicio realizado}

El sistema debe almacenar la información básica de cada servicio o tratamiento realizado a un cliente del salón.

Para cada servicio realizado el sistema debe registrar:

\begin{itemize}
    \item Identificador del servicio realizado.
    \item Cliente del salón asociado.
    \item Cliente profesional o peluquería en la que se realiza el servicio.
    \item Usuario que registra o realiza el servicio, cuando corresponda.
    \item Fecha de realización del servicio.
    \item Tipo de servicio o tratamiento aplicado.
    \item Observaciones generales del servicio.
    \item Resultado o descripción final del trabajo realizado.
\end{itemize}

\subsubsection*{Ficha técnica del servicio}

El sistema debe permitir almacenar la información técnica asociada a cada servicio realizado. Esta información recoge los detalles profesionales necesarios para documentar correctamente el tratamiento aplicado.

Para cada ficha técnica el sistema debe registrar:

\begin{itemize}
    \item Identificador de la ficha técnica.
    \item Servicio realizado asociado.
    \item Descripción técnica del tratamiento.
    \item Fórmula o formulación utilizada, en caso de existir.
    \item Observaciones técnicas asociadas al servicio.
    \item Fecha de creación del registro técnico.
    \item Fecha de última actualización de la ficha técnica.
\end{itemize}

\subsubsection*{Producto utilizado en servicio}

El sistema debe permitir registrar los productos utilizados durante la realización de un servicio técnico.

Para cada producto utilizado en un servicio el sistema debe almacenar:

\begin{itemize}
    \item Identificador del registro de uso de producto.
    \item Servicio realizado asociado.
    \item Producto utilizado.
    \item Cantidad utilizada, en caso de registrarse.
    \item Observaciones asociadas al uso del producto.
\end{itemize}

\subsubsection*{Sugerencia de producto}

El sistema debe permitir almacenar sugerencias de productos generadas a partir del historial técnico de un cliente del salón. Estas sugerencias se presentan como apoyo a la consulta profesional y no sustituyen el criterio técnico del usuario.

Para cada sugerencia de producto el sistema debe registrar:

\begin{itemize}
    \item Identificador de la sugerencia.
    \item Cliente del salón asociado a la sugerencia.
    \item Producto sugerido.
    \item Motivo o justificación de la sugerencia.
    \item Fecha de generación de la sugerencia.
\end{itemize}

% ------------------------------------------------------------------------------
\subsection{M6. Promociones, notificaciones y formaciones}
% ------------------------------------------------------------------------------

El módulo de promociones, notificaciones y formaciones gestiona la información relacionada con la comunicación entre el distribuidor y las peluquerías usuarias del sistema. Su objetivo es permitir la difusión de promociones comerciales, la segmentación de clientes mediante categorías, el envío de avisos y recordatorios relevantes y la gestión de actividades formativas dirigidas a los clientes profesionales.

La información gestionada por este módulo permite mostrar promociones activas, asociarlas a categorías de clientes, registrar notificaciones relacionadas con eventos importantes del sistema, programar recordatorios y mantener un catálogo de formaciones disponibles junto con la posible inscripción de usuarios autorizados.

\subsubsection*{Promoción}

El sistema debe almacenar la información básica asociada a cada promoción comercial definida por el distribuidor.

Para cada promoción el sistema debe registrar:

\begin{itemize}
    \item Identificador de la promoción.
    \item Título o nombre de la promoción.
    \item Descripción de la promoción.
    \item Tipo de promoción.
    \item Valor o beneficio promocional asociado.
    \item Fecha de inicio de vigencia.
    \item Fecha de fin de vigencia.
    \item Estado de la promoción.
    \item Producto asociado, cuando corresponda.
    \item Línea de producto asociada, cuando corresponda.
    \item Cliente asociado, cuando corresponda.
    \item Categoría de clientes asociada, cuando corresponda.
    \item Fecha de creación de la promoción.
    \item Fecha de última actualización de la promoción.
\end{itemize}

\subsubsection*{Categoría de cliente}

El sistema debe permitir almacenar categorías de clientes utilizadas para segmentar a las peluquerías usuarias del sistema en función de criterios comerciales o de uso definidos por el distribuidor.

Para cada categoría de cliente el sistema debe registrar:

\begin{itemize}
    \item Identificador de la categoría.
    \item Nombre de la categoría.
    \item Descripción de la categoría, en caso de existir.
    \item Criterios o condiciones asociados a la categoría, cuando corresponda.
    \item Fecha de creación de la categoría.
    \item Fecha de última actualización de la categoría.
\end{itemize}

\subsubsection*{Categoría de cliente}

El sistema debe permitir almacenar categorías de clientes utilizadas para segmentar a las peluquerías usuarias del sistema en función de criterios comerciales o de uso definidos por el distribuidor.

Para cada categoría de cliente el sistema debe registrar:

\begin{itemize}
    \item Identificador de la categoría.
    \item Nombre de la categoría.
    \item Descripción de la categoría, en caso de existir.
    \item Criterios o condiciones asociados a la categoría, cuando corresponda.
    \item Fecha de creación de la categoría.
    \item Fecha de última actualización de la categoría.
\end{itemize}

\subsubsection*{Notificación}

El sistema debe almacenar la información asociada a las notificaciones o avisos generados para los usuarios del sistema.

Para cada notificación el sistema debe registrar:

\begin{itemize}
    \item Identificador de la notificación.
    \item Usuario destinatario de la notificación.
    \item Título de la notificación.
    \item Contenido o mensaje de la notificación.
    \item Tipo de notificación.
    \item Canal de notificación utilizado.
    \item Fecha y hora de creación de la notificación.
    \item Fecha y hora de envío, cuando corresponda.
    \item Estado de lectura de la notificación.
    \item Fecha de lectura, cuando corresponda.
\end{itemize}

\subsubsection*{Recordatorio}

El sistema debe permitir almacenar recordatorios asociados a eventos o acciones relevantes del sistema, como visitas comerciales programadas o fechas de entrega de pedidos.

Para cada recordatorio el sistema debe registrar:

\begin{itemize}
    \item Identificador del recordatorio.
    \item Usuario destinatario del recordatorio.
    \item Título o descripción breve del recordatorio.
    \item Contenido del recordatorio.
    \item Fecha y hora programadas para el recordatorio.
    \item Estado del recordatorio.
    \item Fecha de creación del recordatorio.
\end{itemize}

\subsubsection*{Formación}

El sistema debe almacenar la información básica asociada a cada actividad formativa publicada por el distribuidor.

Para cada formación el sistema debe registrar:

\begin{itemize}
    \item Identificador de la formación.
    \item Título o nombre de la formación.
    \item Descripción de la actividad formativa.
    \item Fecha y hora de celebración.
    \item Ubicación o modalidad de realización.
    \item Contenido o información adicional de la formación.
    \item Estado de la formación.
    \item Fecha de creación del registro de formación.
    \item Fecha de última actualización de la formación.
\end{itemize}

\subsubsection*{Inscripción a formación}

El sistema debe permitir registrar la inscripción de usuarios autorizados a las actividades formativas publicadas por el distribuidor.

Para cada inscripción a formación el sistema debe registrar:

\begin{itemize}
    \item Identificador de la inscripción.
    \item Formación asociada.
    \item Usuario inscrito.
    \item Fecha de inscripción.
    \item Estado de la inscripción.
    \item Observaciones asociadas a la inscripción, cuando corresponda.
\end{itemize}

% ------------------------------------------------------------------------------
\subsection{M7. Administración e integraciones}
% ------------------------------------------------------------------------------

El módulo de administración e integraciones gestiona la información necesaria para la configuración general del sistema, la consulta de información operativa y la relación con herramientas o servicios externos utilizados en la actividad del distribuidor. Su objetivo es facilitar las tareas de administración, permitir la generación de información de apoyo a la toma de decisiones y soportar el intercambio de datos con sistemas externos.

La información gestionada por este módulo permite registrar configuraciones del sistema, mantener información sobre integraciones externas disponibles y generar propuestas de pedido a proveedor a partir de la actividad registrada en la plataforma.

\subsubsection*{Configuración del sistema}

El sistema debe almacenar la información necesaria para mantener determinados parámetros generales de configuración de la plataforma.

Para cada configuración del sistema el sistema debe registrar:

\begin{itemize}
    \item Identificador de la configuración.
    \item Clave o nombre del parámetro de configuración.
    \item Valor asociado a la configuración.
    \item Descripción de la configuración.
    \item Fecha de creación del registro.
    \item Fecha de última actualización de la configuración.
\end{itemize}

\subsubsection*{Integración externa}

El sistema debe almacenar la información asociada a las integraciones disponibles con herramientas o servicios externos utilizados por el distribuidor.

Para cada integración externa el sistema debe registrar:

\begin{itemize}
    \item Identificador de la integración.
    \item Nombre de la integración.
    \item Tipo de integración o servicio externo.
    \item Descripción de la integración.
    \item Estado de la integración.
    \item Información básica de configuración o conexión.
    \item Fecha de creación del registro.
    \item Fecha de última actualización de la integración.
\end{itemize}

\subsubsection*{Exportación o sincronización}

El sistema debe permitir registrar las operaciones de exportación o sincronización de información realizadas con sistemas externos.

Para cada operación de exportación o sincronización el sistema debe registrar:

\begin{itemize}
    \item Identificador de la operación.
    \item Integración externa asociada.
    \item Tipo de operación realizada.
    \item Tipo de información intercambiada.
    \item Fecha y hora de ejecución.
    \item Estado de la operación.
    \item Resultado o mensaje asociado a la operación.
\end{itemize}

\subsubsection*{Propuesta de pedido a proveedor}

El sistema debe permitir almacenar propuestas de pedido generadas para facilitar la reposición de productos al proveedor, basadas en la información comercial registrada en el sistema o introducidas manualmente por el distribuidor.

Estas propuestas deben reflejar la estructura habitual de las hojas de pedido utilizadas en la operativa real del distribuidor.

Para cada propuesta de pedido a proveedor el sistema debe registrar:

\begin{itemize}

    \item Identificador de la propuesta.
    \item Fecha de generación de la propuesta.
    \item Usuario que genera la propuesta, cuando corresponda.
    \item Estado de la propuesta.
    \item Importe total estimado del pedido.
    \item Número total de unidades incluidas en la propuesta.
    \item Observaciones generales asociadas a la propuesta.
    \item Fecha de última actualización de la propuesta.

\end{itemize}

\subsubsection*{Línea de propuesta de pedido a proveedor}

El sistema debe permitir almacenar el detalle de productos incluidos en cada propuesta de pedido a proveedor, siguiendo la estructura de las hojas de pedido utilizadas por el distribuidor.

Para cada línea de propuesta el sistema debe registrar:

\begin{itemize}

    \item Identificador de la línea de propuesta.
    \item Propuesta de pedido a proveedor asociada.
    \item Producto incluido en la propuesta.
    \item Referencia del producto.
    \item Descripción del producto.
    \item Categoría o tipo de producto (técnico, reventa, marketing, etc.).
    \item Cantidad propuesta del producto.
    \item Precio unitario del producto.
    \item Importe total de la línea.
    \item Motivo o criterio de reposición utilizado (consumo histórico, rotación, reposición manual, etc.).

\end{itemize}

% ------------------------------------------------------------------------------
\subsection{M8. Funcionalidades adicionales}
% ------------------------------------------------------------------------------

Este módulo recoge un conjunto de funcionalidades identificadas durante el proceso de análisis de requisitos que, por su carácter complementario, no forman parte de ninguno de los módulos principales del sistema. Su objetivo es ampliar el valor de la aplicación mediante herramientas avanzadas de apoyo al trabajo del profesional y nuevas funcionalidades orientadas tanto a la atención al cliente como a la operativa comercial.

La información asociada a este módulo podrá utilizarse para apoyar la recomendación de tratamientos, la simulación visual de resultados, el reconocimiento de productos mediante cámara y la reserva online de citas por parte de clientes finales. Dado su carácter evolutivo, estas funcionalidades podrán incorporarse en versiones posteriores del sistema.

\subsubsection*{Consulta de recomendación de tratamientos}

El sistema podrá permitir registrar consultas de recomendación de tratamientos realizadas por los profesionales a partir de determinadas características técnicas del cabello del cliente.

Para cada consulta de recomendación el sistema podrá registrar:

\begin{itemize}
    \item Identificador de la consulta.
    \item Usuario profesional que realiza la consulta.
    \item Cliente del salón asociado, cuando corresponda.
    \item Tipo de cabello.
    \item Estado o condición del cabello.
    \item Tratamientos previos relevantes, cuando se indiquen.
    \item Observaciones o parámetros técnicos adicionales, cuando correspondan.
    \item Fecha de realización de la consulta.
\end{itemize}

\subsubsection*{Recomendación de tratamiento}

El sistema podrá almacenar las recomendaciones generadas a partir de una consulta de recomendación de tratamientos. Estas recomendaciones podrán incluir productos o combinaciones de tratamientos disponibles en el catálogo.

Para cada recomendación de tratamiento el sistema podrá registrar:

\begin{itemize}
    \item Identificador de la recomendación.
    \item Consulta de recomendación asociada.
    \item Producto recomendado, cuando corresponda.
    \item Descripción o justificación de la recomendación.
    \item Nivel de relevancia o prioridad de la sugerencia, cuando corresponda.
    \item Fecha de generación de la recomendación.
\end{itemize}

\subsubsection*{Simulación de coloración}

El sistema podrá permitir registrar simulaciones visuales realizadas para previsualizar posibles cambios coloración capilar.

Para cada simulación el sistema podrá registrar:

\begin{itemize}
    \item Identificador de la simulación.
    \item Imagen utilizada para la simulación, cuando se conserve.
    \item Descripción del cambio simulado.
    \item Resultado generado o referencia al resultado visual obtenido.
    \item Fecha de realización de la simulación.
\end{itemize}

\subsubsection*{Reconocimiento de producto mediante cámara}

El sistema podrá permitir registrar operaciones de reconocimiento de productos realizadas a partir de imágenes o códigos de barras capturados con la cámara del dispositivo.

Para cada operación de reconocimiento el sistema podrá registrar:

\begin{itemize}
    \item Identificador de la operación de reconocimiento.
    \item Usuario que realiza la operación.
    \item Imagen capturada o referencia al recurso utilizado, cuando corresponda.
    \item Código de barras detectado, cuando corresponda.
    \item Producto identificado, cuando se obtenga coincidencia.
    \item Finalidad de la operación, como consulta o apoyo a pedido.
    \item Fecha y hora de la operación.
\end{itemize}

\subsubsection*{Reserva online}

El sistema podrá permitir registrar solicitudes de cita realizadas por clientes finales de una peluquería usuaria del sistema.

Para cada reserva online el sistema podrá registrar:

\begin{itemize}
    \item Identificador de la reserva.
    \item Cliente profesional o peluquería a la que se dirige la reserva.
    \item Cliente final que solicita la cita.
    \item Nombre de la persona que realiza la reserva.
    \item Teléfono o correo electrónico de contacto.
    \item Fecha y hora solicitadas para la cita.
    \item Servicio solicitado, cuando corresponda.
    \item Observaciones asociadas a la reserva, cuando existan.
    \item Estado de la reserva.
    \item Fecha de creación de la solicitud.
\end{itemize}